% This must be in the first 5 lines to tell arXiv to use pdfLaTeX, which is strongly recommended.
\pdfoutput=1
% In particular, the hyperref package requires pdfLaTeX in order to break URLs across lines.

\documentclass[11pt]{article}

% Change "review" to "final" to generate the final (sometimes called camera-ready) version.
% Change to "preprint" to generate a non-anonymous version with page numbers.
\usepackage{acl}

% Standard package includes
\usepackage{times}
\usepackage{latexsym}
% \usepackage[pagebackref,breaklinks,colorlinks,citecolor=cvprblue]{hyperref}
\usepackage{multirow}
\usepackage{booktabs}
% For proper rendering and hyphenation of words containing Latin characters (including in bib files)
\usepackage[T1]{fontenc}
% For Vietnamese characters
% \usepackage[T5]{fontenc}
% See https://www.latex-project.org/help/documentation/encguide.pdf for other character sets

% This assumes your files are encoded as UTF8
\usepackage[utf8]{inputenc}

\usepackage{microtype}
\usepackage{graphicx}
\usepackage{subcaption}
\usepackage{amsmath}
\usepackage{cleveref}
\usepackage{tcolorbox}

\usepackage{amssymb}
\usepackage{xcolor}
\usepackage{pifont}
\definecolor{my_green}{RGB}{51,102,0}
\definecolor{my_red}{RGB}{204, 0, 0}
\newcommand{\cmark}{\textcolor{my_green}{\ding{51}}} % ✔
\newcommand{\xmark}{\textcolor{my_red}{\ding{55}}} % ✘
\newcommand{\lilei}[1]{\textcolor{violet}{\bf \small [#1 --lilei]}}
\newcommand{\lpk}[1]{\textcolor{blue}{\bf \small [#1 --lpk]}}

% This is also not strictly necessary, and may be commented out.
% However, it will improve the aesthetics of text in
% the typewriter font.
\usepackage{inconsolata}

% If the title and author information does not fit in the area allocated, uncomment the following
%
%\setlength\titlebox{<dim>}
%
% and set <dim> to something 5cm or larger.

\title{Multimodal ArXiv: A Dataset for Improving Scientific Comprehension of Large Vision-Language Models}

% Author information can be set in various styles:
% For several authors from the same institution:
% \author{Author 1 \and ... \and Author n \\
%         Address line \\ ... \\ Address line}
% if the names do not fit well on one line use
%         Author 1 \\ {\bf Author 2} \\ ... \\ {\bf Author n} \\
% For authors from different institutions:
% \author{Author 1 \\ Address line \\  ... \\ Address line
%         \And  ... \And
%         Author n \\ Address line \\ ... \\ Address line}
% To start a separate ``row'' of authors use \AND, as in
% \author{Author 1 \\ Address line \\  ... \\ Address line
%         \AND
%         Author 2 \\ Address line \\ ... \\ Address line \And
%         Author 3 \\ Address line \\ ... \\ Address line}

\author{
Lei Li\thanks{Equal Contribution.}$^\dagger$,  Yuqi Wang$^*$$^\dagger$, Runxin Xu$^\ddag$, Peiyi Wang$^\ddag$\\\textbf{Xiachong Feng$^\dagger$, Lingpeng Kong$^\dagger$, Qi Liu$^\dagger$}\\ 
$^\dagger$The University of Hong Kong \\
$^\dagger$Peking University\\
 \texttt{\{nlp.lilei, runxinxu, wangpeiyi9979, xiachongfeng1996\}@gmail.com} \\ 
 \texttt{wangyuqi@connect.hku.hk} \quad \texttt{\{lpk, liuqi\}@cs.hku.hk}\\
  }
\def\DatasetName{ArXivCap\ }

\begin{document}
\maketitle
\input{00_abstract}



\input{01_intro}

\input{02_related}
\input{dataset}

\input{04_experiment}

\input{10_conclusion}

% Bibliography entries for the entire Anthology, followed by custom entries
%\bibliography{anthology,custom}
% Custom bibliography entries only
\bibliography{11_references}
\appendix

\section{Details of Multimodal ArXiv}
\label{apx:arxiv_cap}


\subsection{Caption Cleaning}
\label{apx:caption_clean}

We apply a Python tool to clean the original caption and Table~\ref{tab:pylatexenc_clean} illustrates the caption before and after cleaning. 
\input{tables/pylatexenc_clean}

\subsection{Caption Word Cloud}
\label{apx:caption_word_cloud}

We visualize the word cloud of captions in our ArXivCap dataset in Figure~\ref{fig:word_cloud}, it can be seen that the captions have a diverse vocabulary for describing the different figures in the academic papers.
\begin{figure}[t!]
    \centering
    \includegraphics[width=\linewidth]{figs/word_cloud.jpg}
    \caption{Word cloud visualization of captions in our ArXivCap dataset.}
    \label{fig:word_cloud}
\end{figure}


\subsection{Illustration Cases from ArXivCap and ArXivQA}
\label{apx:case_illustrations}
We provide illustrated cases from our dataset for a better understanding.
Figure~\ref{fig:single_figure_pair} demonstrates a typical single figure-caption pair, and Figure~\ref{fig:multi_figure_case} shows the multiple figure-caption case.
Figure~\ref{fig:qa_case_1}, \ref{fig:qa_case_2}, \ref{fig:qa_case_3} and \ref{fig:qa_case_4} illustrate the cases from our ArXivQA dataset, covering different figure types and containing diverse questions.



\begin{figure}[htb!]
    \centering
    \includegraphics[width=\linewidth]{figs/single-figure.pdf}
    \caption{A single-figure caption pair in our ArXivCap dataset.}
\label{fig:single_figure_pair}
\end{figure}


\begin{figure*}[t!]
    \centering
    \includegraphics[width=\linewidth]{figs/multiple-figure.pdf}
    \caption{A multiple-figure caption pair in our ArXivCap dataset.}
    \label{fig:multi_figure_case}
\end{figure*}


\begin{figure*}[tbh!]
    \centering
    \includegraphics[width=\linewidth]{figs/arxivqa_case1.pdf}
    \caption{A case from ArXivQA.}
    \label{fig:qa_case_1}
\end{figure*}



\begin{figure*}[tbh!]
    \centering
    \includegraphics[width=\linewidth]{figs/arxivqa_case2.pdf}
    \caption{A case from ArXivQA.}
    \label{fig:qa_case_2}
\end{figure*}



\begin{figure*}[tbh!]
    \centering
    \includegraphics[width=\linewidth]{figs/arxivqa_case3.pdf}
    \caption{A case from ArXivQA.}
    \label{fig:qa_case_3}
\end{figure*}



\begin{figure*}[tbh!]
    \centering
    \includegraphics[width=\linewidth]{figs/arxivqa_case4.pdf}
    \caption{A case from ArXivQA.}
    \label{fig:qa_case_4}
\end{figure*}
\input{tables/domain-full-name}


\input{tables/arxivqa_template}
\subsection{ArXivQA Prompting Template}
\label{apx:prompt_template}
The prompt used to query GPT-4V is provided in Table~\ref{tab:prompt_for_ArXivqa}.

\section{Evaluation Details}
\label{apx:evaluation_details}
\subsection{Details of Evaluated Models}


\noindent\textbf{BLIP2}~\citep{li2023blip2}, 
introduces a lightweight Q-Former designed to bridge modality gaps and leverages frozen LLMs. Leveraging LLMs, BLIP-2 can conduct zero-shot image-to-text generation using natural language prompts. We select the \texttt{BLIP2-OPT-6.7B} version for evaluation.\footnote{\url{https://huggingface.co/Salesforce/blip2-opt-6.7b}}


\noindent\textbf{InstructBLIP}~\citep{dai2023instructblip} employs an instruction-aware visual feature extraction module based on BLIP2~\citep{li2023blip2} and is trained with unified multimodal instruction tuning datasets. We choose \texttt{InstructBLIP-Vicuna-7B} for evaluation.\footnote{\url{https://huggingface.co/Salesforce/instructblip-vicuna-7b}}.



\noindent\textbf{LLaVA}~\citep{liu2023llava},
adopts Vicuna models as the backbone LLM and is trained on the ChatGPT/GPT-4 generated instruction tuning dataset.
LLaVA-v1.5~\citep{liu2023llava15} improves on LLaVA models by employing curated task datasets and an enhanced modality alignment module.
We evaluate both \texttt{LLaVA-v1.5-7B}\footnote{\url{https://huggingface.co/liuhaotian/llava-v1.5-7b}} and \texttt{LLaVA-v1.5-13B}.\footnote{\url{https://huggingface.co/liuhaotian/llava-v1.5-13b}}


\noindent\textbf{Flamingo}~\citep{Alayrac2022FlamingoAV} pioneers the development of LVLMs by introducing a cross-gated layer for LLMs to produce visual-grounded text. The training dataset
consists of interleaved visual data and text from the web pages, enabling it to generate free-form text as the output. We select the open-source implementation \texttt{OpenFlamingo-9B}~\citep{awadalla2023openflamingo} for evaluation.\footnote{\url{https://huggingface.co/openflamingo/OpenFlamingo-9B-vitl-mpt7b}}

\noindent\textbf{IDEFICS} is another open-sourced implementation of Flamingo~\citep{Alayrac2022FlamingoAV} Trained on publicly available image-text alignment pairs and instruction tuning datasets, it demonstrates comparable results with the original closed-source model on various image-text benchmarks. We select the \texttt{IDEFICS-Instruct-9B} for evaluation.\footnote{\url{https://huggingface.co/HuggingFaceM4/idefics-9b-instruct}}.

\noindent\textbf{Qwen-VL-Chat}~\citep{Qwen-VL} is
a bilingual LVLM that supports both English and Chinese built on the Qwen LLM~\citep{qwen}. 
During the training phase, Qwen-VL-Chat adopts a packing strategy to create multiple images as inputs,
improving its ability to understand the vision context. We select the \texttt{Qwen-VL-Chat-7B} for evaluation.\footnote{\url{https://github.com/QwenLM/Qwen-VL}}



\noindent\textbf{GPT-4V}
~\citep{gpt4v}, the proprietary vision-language models developed by OpenAI, which are shown to be powerful on various multi-modal tasks~\citep{yang2023dawn}. The API version we queried is \texttt{gpt-4-vision-preview}.

\noindent\textbf{Bard}~\citep{bard} (now Gemini), a commercial LVLM developed by Google. We utilize the unofficial API\footnote{\url{https://github.com/dsdanielpark/Bard-API}} querying the model with our task prompts, accessed on \texttt{2023-11-17}.


\subsection{Task Prompts}
We evaluate all the models with the same task prompts in our experiments, and the prompts for our four tasks are listed below:

\begin{itemize}
    \item Single-Figure Captioning:  \texttt{Create a caption for the provided figure.}
    \item Multiple-Figure Captioning: \texttt{Create a caption for the provided figures.}
    \item Contextualized Captioning: We reuse the prompts in previous captioning tasks depending on the current figure type.
    \item Title Generation: \texttt{
According to the figures and captions, generate a title for this paper. Title:} 
\end{itemize}





\section{Error Analysis}
\label{apx:case}

\subsection{Caption Type Illustration} 
\label{apx:caption_type}
We illustrate captions of four types in our main paper in Figure~\ref{fig:caption_type}. The \emph{Acceptable} caption provides a comprehensive description of the figure presented. The \emph{oversimplified} caption is too short compared with the original human-written caption. 
Furthermore, as show in the third block in Figure~\ref{fig:caption_type}, \emph{Contextual Misinterpretation} refers to captions with unmentioned content in the figure, such as the dataset colored in red.
\emph{Recognition Error} denotes the model wrongly identified the number or text in the figure, such as the misidentified model name in the last block of Figure~\ref{fig:caption_type}.

\subsection{Failure Sample of MathVista}
\label{apx:failure_mathvista}
Figure~\ref{fig:geometry_fail} shows a challenging geometry mathematic reasoning problem where both models fail to produce the correct answer.
Echoing the quantitative results in our main paper, we believe future studies can incorporate more focused corpus for enhancing the geometry and mathematical reasoning ability of LVLMs.
\begin{figure*}[t!]
    \centering
    \includegraphics[width=\linewidth]{figs/caption_error_type.pdf}
    \caption{Illustration of acceptable and three error types of generated captions.}
    \label{fig:caption_type}
\end{figure*}


\begin{figure}
    \centering
    \includegraphics[width=\linewidth]{figs/mathvista_both_wrong.pdf}
    \caption{A failure case on the geometry problem-solving task.}
\label{fig:geometry_fail}
\end{figure}
\end{document}
